%% toc: P1 6 6.1,6.3,6.4
\section{Semantic data, models, and undischarged assumptions}
\label{sec:p01-w02-p01-016-semantic-data}

The definitions of \cref{sec:levels-03} rest on three notions that they use
without defining. Semantic consequence quantifies over ``every complete tuple
of semantic data permitted by that semantics'' and asks whether each tuple
\emph{satisfies} the premises; derivability counts the \emph{undischarged
assumptions} of a derivation. The definition of a semantics in
\cref{sec:levels-01} names the same tuples, and \cref{sec:levels-04} speaks of
models, but neither says what a datum is, when it is complete, what a
variable assignment ranges over, when a set of formulas is satisfiable, or
how an assumption comes to be discharged. This section supplies those
definitions for the three semantics named in \cref{sec:levels-03}
(propositional valuations, first-order structures with a variable
assignment, Kripke structures with a world) and for the natural-deduction
calculi \(\mathsf{ND}_{\mathsf{CPL}}\) and \(\mathsf{ND}_{\mathsf{FOL}}\).

The inputs are the inherited definitions of valuations, structures,
assignments, satisfaction, models and natural-deduction derivations
\cite[chs.~``Levels, Objects, and Consequence'', ``Propositional Syntax and
Semantics'', ``Natural Deduction and Sequents'', ``Structures, Satisfaction,
and Translation'', ``Theories and Axiomatization'', ``First-Order
Deduction'']{FrozenFormalLogicSemanticData}. The conclusions are: an exact
account of semantic data and of satisfaction
(\cref{def:p01-w02-p01-016-datum,def:p01-w02-p01-016-satisfaction}); the
structure of assignments and their variants
(\cref{thm:p01-w02-p01-016-variants,thm:p01-w02-p01-016-quantifier-variants});
a proof that a tuple must be complete for consequence to be the intended
relation (\cref{thm:p01-w02-p01-016-local-global}); satisfiability and its
relation to consequence and to models
(\cref{thm:p01-w02-p01-016-unsat-reduction,thm:p01-w02-p01-016-sentences}); and
an account of open and discharged assumptions with the discharge theorem for
natural deduction
(\cref{thm:p01-w02-p01-016-open-bounds,thm:p01-w02-p01-016-discharge}). The
results and proofs are editorial completions built from the inherited
definitions; they are not proofs attributed to an external presentation.

\subsection{Semantic data and satisfaction}
\label{sec:p01-w02-p01-016-data}

A semantics fixes what a formula is evaluated \emph{at}. The thing it is
evaluated at is a tuple, and the tuple has to supply every component that a
truth condition consults; otherwise the question whether the formula holds
has no answer.

\begin{hbdefinition}{def:p01-w02-p01-016-datum}{Semantic datum and complete tuple}
Let \(\mathcal S\) be a semantics for a formal language with formula set
\(F\), in the sense of \cref{sec:levels-01}. The semantics declares a class
\(\mathcal D_{\mathcal S}\) of \emph{permitted semantic data}. Each member is
a tuple whose components have declared types, and the tuple is
\emph{complete}: it supplies, with its declared type, every component that
some satisfaction clause of \(\mathcal S\) consults. A complete tuple of
semantic data for \(\mathcal S\) is exactly a member of
\(\mathcal D_{\mathcal S}\). A tuple that lacks a consulted component, or
supplies one of the wrong type, is not a datum of \(\mathcal S\), and
\(\mathcal S\) assigns it no satisfaction relation. The three semantics named
in \cref{sec:levels-03} declare the following data.
\begin{enumerate}
  \item \emph{Classical propositional logic.} A datum is a single valuation
        \(v\colon\mathsf{Prop}\to\mathbb B\), where
        \(\mathbb B=\{\mathbf F,\mathbf T\}\). Completeness means that
        \(v\) is total: it assigns a value to every propositional variable,
        not only to those occurring in the formulas at hand
        \cite[ch.~``Propositional Syntax and Semantics'', definition ``Truth
        values and valuation'']{FrozenFormalLogicSemanticData}.
  \item \emph{First-order logic over a signature \(\Lang\).} A datum is a
        pair \((\mathcal M,s)\) of an \(\Lang\)-structure \(\mathcal M\),
        whose domain \(\dom\mathcal M\) is nonempty, and a variable
        assignment \(s\) in \(\mathcal M\)
        (\cref{def:p01-w02-p01-016-assignment}). The structure alone is not a
        datum
        \cite[ch.~``Structures, Satisfaction, and Translation'', definitions
        ``First-order structure'' and ``Variable assignment and
        update'']{FrozenFormalLogicSemanticData}.
  \item \emph{Normal modal logic over a class \(\mathcal C\) of frames.} A
        datum is a pair \((\mathcal M,w)\) of a Kripke structure
        \(\mathcal M=(W,R,V)\) with \(W\ne\varnothing\),
        \(R\subseteq W\times W\), \(V\colon\mathsf{Prop}\to\mathcal P(W)\)
        and \((W,R)\in\mathcal C\), together with a world \(w\in W\)
        \cite[ch.~``Modal, Temporal, and Epistemic Logic'', definition
        ``Kripke frame and Kripke structure'']{FrozenFormalLogicSemanticData}.
        The Kripke structure alone is not a datum.
\end{enumerate}
\index{semantic data}\index{complete tuple of semantic data}
\end{hbdefinition}

\begin{hbdefinition}{def:p01-w02-p01-016-assignment}{Variable assignment and assignment variant}
Let \(\mathcal M\) be a first-order structure and let \(\mathsf{Var}\) be the
set of individual variables. A \emph{variable assignment} in \(\mathcal M\) is
a total function
\[
  s\colon\mathsf{Var}\to\dom\mathcal M .
\]
Its domain is the whole of \(\mathsf{Var}\), not only the variables free in
some formula; its codomain is the domain of the structure. For
\(x\in\mathsf{Var}\) and \(a\in\dom\mathcal M\), the \emph{update}
\(s[x\mapsto a]\) is the assignment that sends \(x\) to \(a\) and agrees with
\(s\) on every other variable
\cite[ch.~``Structures, Satisfaction, and Translation'', definition
``Variable assignment and update'']{FrozenFormalLogicSemanticData}. An
assignment \(s'\) in \(\mathcal M\) is an \emph{\(x\)-variant} of \(s\) iff
\(s'(y)=s(y)\) for every variable \(y\ne x\). Every assignment is an
\(x\)-variant of itself.
\index{variable assignment}\index{assignment!variant}
\end{hbdefinition}

\begin{hblemma}{thm:p01-w02-p01-016-variants}{proved=annex:p01-w02-p01-016-variants}{Updates and variants}
Let \(\mathcal M\) be a structure, \(s\) an assignment in \(\mathcal M\),
\(x,y\) variables and \(a,b\in\dom\mathcal M\).
\begin{enumerate}
  \item At least one assignment in \(\mathcal M\) exists.
  \item \(s[x\mapsto a]\) is an \(x\)-variant of \(s\), and every
        \(x\)-variant \(s'\) of \(s\) equals \(s[x\mapsto s'(x)]\). The map
        \(a\mapsto s[x\mapsto a]\) is a bijection from \(\dom\mathcal M\) onto
        the set of \(x\)-variants of \(s\).
  \item \(s[x\mapsto a][x\mapsto b]=s[x\mapsto b]\) and
        \(s[x\mapsto s(x)]=s\).
  \item If \(x\ne y\), then
        \(s[x\mapsto a][y\mapsto b]=s[y\mapsto b][x\mapsto a]\).
\end{enumerate}
\end{hblemma}

Satisfaction is the relation that answers, for each datum and each formula,
whether the formula holds at that datum. It is defined by structural
recursion on formulas, one clause for each formation rule; only the atomic
clauses and the clauses for variable-binding or world-shifting operators
consult components of the datum.

\begin{hbdefinition}{def:p01-w02-p01-016-satisfaction}{Satisfaction relation}
A \emph{satisfaction relation} for a semantics \(\mathcal S\) is a relation
between \(\mathcal D_{\mathcal S}\) and the formula set \(F\); write
\(d\entails_{\mathcal S}\varphi\), read ``\(d\) satisfies \(\varphi\)'', when
it holds, and \(d\not\entails_{\mathcal S}\varphi\) otherwise. For
\(\Gamma\subseteq F\), \(d\entails_{\mathcal S}\Gamma\) means
\(d\entails_{\mathcal S}\gamma\) for every \(\gamma\in\Gamma\); in particular
every datum satisfies \(\varnothing\). The subscript is dropped when one
semantics is fixed. For the three semantics of
\cref{def:p01-w02-p01-016-datum}:
\begin{enumerate}
  \item \(v\entails\varphi\) iff \(\widehat v(\varphi)=\mathbf T\), where
        \(\widehat v\) is the unique extension of \(v\) by the recursive
        truth clauses of \cref{sec:prop-syntax-semantics-02};
  \item \(\mathcal M,s\entails\varphi\) is defined by the clauses of
        \cref{sec:fol-semantics-02}. The clauses that consult the datum are
        the atomic ones,
        \[
          \begin{aligned}
          \mathcal M,s\entails R(t_1,\ldots,t_n)
          &\iff
          (\sem{t_1}^{\mathcal M}_s,\ldots,\sem{t_n}^{\mathcal M}_s)\in
          R^{\mathcal M},\\
          \mathcal M,s\entails r=t
          &\iff
          \sem{r}^{\mathcal M}_s=\sem{t}^{\mathcal M}_s,
          \end{aligned}
        \]
        where \(\sem{t}^{\mathcal M}_s\) is the value of the term \(t\) in
        \(\mathcal M\) under \(s\), and the quantifier clauses,
        \[
          \mathcal M,s\entails\forall x\,\varphi
          \iff
          \mathcal M,s[x\mapsto a]\entails\varphi
          \text{ for every }a\in\dom\mathcal M,
        \]
        with ``some'' in place of ``every'' for \(\exists x\,\varphi\); the
        clauses for \(\bottom,\neg,\land,\lor,\to\) are the Boolean ones,
        with \(\mathcal M,s\entails\neg\varphi\) iff
        \(\mathcal M,s\not\entails\varphi\);
  \item \(\mathcal M,w\entails A\) is classical for the propositional
        connectives at the world \(w\), and
        \(\mathcal M,w\entails\Box A\) iff \(\mathcal M,v\entails A\) for
        every \(v\) with \(wRv\) \cite[ch.~``Modal, Temporal, and Epistemic
        Logic'', definition ``Kripke frame and Kripke
        structure'']{FrozenFormalLogicSemanticData}. The inherited definition
        does not display the atomic clause; here it is read as
        \(\mathcal M,w\entails p\) iff \(w\in V(p)\), the clause that the
        inherited intuitionistic semantics states for a valuation of the same
        type \(\mathsf{Prop}\to\mathcal P(W)\) \cite[ch.~``Intuitionistic
        Logic'', forcing relation following definition ``Intuitionistic
        Kripke frame and structure'']{FrozenFormalLogicSemanticData}.
\end{enumerate}
In each case the semantic consequence relation of \cref{sec:levels-03} is
\[
  \Gamma\entails_{\mathcal S}\varphi
  \iff
  \text{for every }d\in\mathcal D_{\mathcal S}:\;
  d\entails_{\mathcal S}\Gamma\text{ implies }d\entails_{\mathcal S}\varphi .
\]
\index{satisfaction!relation of a semantics}
\end{hbdefinition}

The quantifier clause is often stated with variants instead of updates. The
two readings agree.

\begin{hbtheorem}{thm:p01-w02-p01-016-quantifier-variants}{proved=annex:p01-w02-p01-016-quantifier-variants}{Quantifier clauses through variants}
For every structure \(\mathcal M\), assignment \(s\), variable \(x\) and
formula \(\varphi\):
\begin{enumerate}
  \item \(\mathcal M,s\entails\forall x\,\varphi\) iff
        \(\mathcal M,s'\entails\varphi\) for every \(x\)-variant \(s'\) of
        \(s\);
  \item \(\mathcal M,s\entails\exists x\,\varphi\) iff
        \(\mathcal M,s'\entails\varphi\) for some \(x\)-variant \(s'\) of
        \(s\).
\end{enumerate}
\end{hbtheorem}

Completeness of a tuple is a requirement on the definition, not on what has
to be inspected in a particular case. The coincidence lemmas
(\cref{lem:prop-coincidence,lem:fol-coincidence}) show that satisfaction of
a formula depends only on the components it mentions, and this makes finite
checks possible for finitely many propositional formulas.

\begin{hbtheorem}{thm:p01-w02-p01-016-finite-check}{proved=annex:p01-w02-p01-016-finite-check}{Finitely many valuations suffice for propositional consequence}
Let \(\Gamma\cup\{\varphi\}\) be a set of propositional formulas and let
\(P\subseteq\mathsf{Prop}\) be a finite set containing every propositional
variable that occurs in a member of \(\Gamma\cup\{\varphi\}\). For
\(u\colon P\to\mathbb B\), let \(u^{+}\) be the valuation that agrees with
\(u\) on \(P\) and sends every variable outside \(P\) to \(\mathbf F\). Then
\(\Gamma\entails\varphi\) iff, for each of the \(2^{\lvert P\rvert}\)
functions \(u\colon P\to\mathbb B\), \(u^{+}\entails\Gamma\) implies
\(u^{+}\entails\varphi\).
\end{hbtheorem}

The truth table in \cref{sec:prop-syntax-semantics-02} is this check with
\(\Gamma=\varnothing\) and \(P=\{P,Q\}\). Partial data serve only because the
coincidence lemma relates them to complete ones; they are not themselves data
of the semantics.

The next result shows what goes wrong when the components of a datum are not
held fixed together. Quantifying first over structures, and only inside that
over assignments, yields a different relation.

\begin{hbtheorem}{thm:p01-w02-p01-016-local-global}{proved=annex:p01-w02-p01-016-local-global}{Consequence must quantify over whole tuples}
\leavevmode
\begin{enumerate}
  \item Let \(R\) be a unary relation symbol and \(x\) a variable. For every
        structure \(\mathcal M\): if \(\mathcal M,s\entails R(x)\) for every
        assignment \(s\), then \(\mathcal M,s\entails\forall x\,R(x)\) for
        every assignment \(s\). Nevertheless
        \(\{R(x)\}\not\entails\forall x\,R(x)\).
  \item Let \(p\in\mathsf{Prop}\) and let \(\mathcal C\) be the class of all
        frames. For every Kripke structure \(\mathcal M\) on a frame in
        \(\mathcal C\): if \(\mathcal M,w\entails p\) for every world \(w\),
        then \(\mathcal M,w\entails\Box p\) for every world \(w\).
        Nevertheless \(\{p\}\not\entails^{\ell}_{\mathcal C}\Box p\), where
        \(\entails^{\ell}_{\mathcal C}\) is the consequence relation of
        \cref{sec:levels-03} for the data \((\mathcal M,w)\).
\end{enumerate}
\end{hbtheorem}

The relation in which the whole tuple is quantified at once is the one
defined in \cref{sec:levels-03}. For modal data it is the local consequence
\(\entails^{\ell}_{\mathcal C}\) of \cref{sec:modal-temporal-epistemic-01}; the
structure-wise relation is the global consequence
\(\entails^{g}_{\mathcal C}\), and the inherited text records that the two
can differ \cite[ch.~``Modal, Temporal, and Epistemic Logic'', definition
``Local and global modal consequence'']{FrozenFormalLogicSemanticData}.
Clause 2 of \cref{thm:p01-w02-p01-016-local-global} is a witness.

\begin{hbwarning}{def:p01-w02-p01-016-incomplete-tuple}{A structure without an assignment does not satisfy an open formula}
Let \(\mathcal M\) have domain \(\{0,1\}\) and \(R^{\mathcal M}=\{0\}\), and
let \(s(x)=0\) and \(s'(x)=1\). Then \(\mathcal M,s\entails R(x)\) and
\(\mathcal M,s'\not\entails R(x)\). The question whether ``\(\mathcal M\)
satisfies \(R(x)\)'' therefore has no answer: \((\mathcal M)\) is not a
complete tuple. Writing \(\mathcal M\entails R(x)\) and meaning ``for some
assignment'' or ``for every assignment'' silently changes the consequence
relation, as \cref{thm:p01-w02-p01-016-local-global} shows. Only for a
sentence is the choice of assignment immaterial, by the assignment
independence corollary in \cref{sec:fol-semantics-02}.
\end{hbwarning}

\subsection{Models, satisfiability, validity, and consequence}
\label{sec:p01-w02-p01-016-models}

\begin{hbdefinition}{def:p01-w02-p01-016-satisfiable}{Satisfiability}
Let \(\mathcal S\) be a semantics with data \(\mathcal D_{\mathcal S}\). A set
\(\Gamma\) of formulas is \emph{satisfiable} in \(\mathcal S\) iff
\(d\entails_{\mathcal S}\Gamma\) for some \(d\in\mathcal D_{\mathcal S}\),
and \emph{unsatisfiable} otherwise. A formula \(\varphi\) is satisfiable iff
\(\{\varphi\}\) is. For the three semantics of
\cref{def:p01-w02-p01-016-datum} this reads: some valuation satisfies
\(\Gamma\) \cite[ch.~``Propositional Syntax and Semantics'', definition
``Semantic notions'']{FrozenFormalLogicSemanticData}; some structure and
assignment satisfy every member of \(\Gamma\) \cite[ch.~``Theories and
Axiomatization'', definition ``Theory, model, satisfiability, and structure
theory'']{FrozenFormalLogicSemanticData}; some Kripke structure on a frame
in \(\mathcal C\) and some world of it satisfy every member of \(\Gamma\).
\index{satisfiable}\index{unsatisfiable}
\end{hbdefinition}

Validity is the case \(\Gamma=\varnothing\) of consequence
(\cref{sec:levels-03}): \(\varphi\) is valid iff every datum satisfies it.
A \emph{model} is a first-order notion tied to sentences: a structure
\(\mathcal M\) is a model of a set \(T\) of sentences iff
\(\mathcal M\entails\sigma\) for every \(\sigma\in T\), where
\(\mathcal M\entails\sigma\) means \(\mathcal M,s\entails\sigma\) for every
assignment \(s\) (\cref{sec:levels-04}, \cref{def:fol-theory-satisfiability}).
A structure is not called a model without a stated sentence or theory, and a
datum of the propositional or modal semantics is not called a model in this
handbook's registered vocabulary.

\begin{hbtheorem}{thm:p01-w02-p01-016-unsat-reduction}{proved=annex:p01-w02-p01-016-unsat-reduction}{Consequence, validity and satisfiability}
Let \(\mathcal S\) be a semantics whose satisfaction relation obeys the
classical negation clause, \(d\entails_{\mathcal S}\neg\varphi\) iff
\(d\not\entails_{\mathcal S}\varphi\), for every datum \(d\) and formula
\(\varphi\). This holds for the three semantics of
\cref{def:p01-w02-p01-016-datum}. For all \(\Gamma,\Delta,\varphi\):
\begin{enumerate}
  \item \(\Gamma\entails_{\mathcal S}\varphi\) iff
        \(\Gamma\cup\{\neg\varphi\}\) is unsatisfiable;
  \item \(\varphi\) is valid iff \(\neg\varphi\) is unsatisfiable, and
        \(\varphi\) is satisfiable iff \(\neg\varphi\) is not valid;
  \item if \(\Gamma\subseteq\Delta\) and \(\Gamma\entails_{\mathcal S}\varphi\),
        then \(\Delta\entails_{\mathcal S}\varphi\); if \(\Delta\) is
        satisfiable, so is every subset of \(\Delta\);
  \item if \(\Gamma\) is unsatisfiable, then
        \(\Gamma\entails_{\mathcal S}\psi\) for every formula \(\psi\).
\end{enumerate}
Clauses 3 and 4 hold for every semantics, without the negation clause.
\end{hbtheorem}

\begin{hbtheorem}{thm:p01-w02-p01-016-sentences}{proved=annex:p01-w02-p01-016-sentences}{Satisfiability and consequence for sentences through models}
Let \(T\cup\{\sigma\}\) be a set of sentences of a first-order signature
\(\Lang\).
\begin{enumerate}
  \item \(T\) is satisfiable iff \(T\) has a model.
  \item \(T\entails\sigma\) iff every model of \(T\) is a model of
        \(\sigma\), that is,
        \(\operatorname{Mod}(T)\subseteq\operatorname{Mod}(\{\sigma\})\).
  \item \(\sigma\) is valid iff \(\mathcal M\entails\sigma\) for every
        \(\Lang\)-structure \(\mathcal M\).
\end{enumerate}
\end{hbtheorem}

The step from satisfiability to a model in clause 1 uses the nonemptiness of
domains: a model is a structure in which every assignment satisfies \(T\), and
there must be an assignment to start from (clause 1 of
\cref{thm:p01-w02-p01-016-variants}).

\begin{hbexample}{def:p01-w02-p01-016-worked-models}{A consequence checked through unsatisfiability}
Let \(\Lang\) have unary relation symbols \(P,Q\) and a constant symbol
\(c\), and let \(\Gamma=\{\forall x\,(P(x)\to Q(x)),\,P(c)\}\). By clause 1
of \cref{thm:p01-w02-p01-016-unsat-reduction}, \(\Gamma\entails Q(c)\) holds
iff \(\Gamma\cup\{\neg Q(c)\}\) is unsatisfiable. Suppose
\(\mathcal M,s\) satisfied all three formulas and put \(a=c^{\mathcal M}\).
The universal clause at \(a\) gives
\(\mathcal M,s[x\mapsto a]\entails P(x)\to Q(x)\). Since
\(\mathcal M,s\entails P(c)\), \(a\in P^{\mathcal M}\), so
\(\mathcal M,s[x\mapsto a]\entails P(x)\), hence
\(\mathcal M,s[x\mapsto a]\entails Q(x)\) and \(a\in Q^{\mathcal M}\). Then
\(\mathcal M,s\entails Q(c)\), contradicting
\(\mathcal M,s\entails\neg Q(c)\). So \(\Gamma\entails Q(c)\).

The premises are sentences, so by \cref{thm:p01-w02-p01-016-sentences} the
same fact says that every model of \(\Gamma\) is a model of \(Q(c)\). The set
\(\Gamma\) is satisfiable: the structure with domain \(\{0\}\),
\(P^{\mathcal M}=Q^{\mathcal M}=\{0\}\) and \(c^{\mathcal M}=0\) is a model.
Dropping the universal premise changes the verdict: \(\{P(c),\neg Q(c)\}\) is
satisfied by the structure with domain \(\{0\}\), \(P^{\mathcal M}=\{0\}\),
\(Q^{\mathcal M}=\varnothing\), \(c^{\mathcal M}=0\), so
\(\{P(c)\}\not\entails Q(c)\).
\end{hbexample}

Satisfiability of a set with free variables is a property of pairs, and it is
not preserved by passing to universal closures.

\begin{hbtheorem}{thm:p01-w02-p01-016-closure-gap}{proved=annex:p01-w02-p01-016-closure-gap}{Satisfiable formulas with unsatisfiable closures}
Let \(R\) be a unary relation symbol and \(x\ne y\) variables.
\begin{enumerate}
  \item \(\{R(x),\neg R(y)\}\) is satisfiable.
  \item \(\{\forall x\,R(x),\forall y\,\neg R(y)\}\) is unsatisfiable, so it
        has no model.
  \item No structure \(\mathcal M\) has
        \(\mathcal M,s\entails R(x)\) and \(\mathcal M,s\entails\neg R(y)\)
        for every assignment \(s\).
\end{enumerate}
\end{hbtheorem}

The set \(\{R(x),\neg R(y)\}\) has no model in the registered sense, because
model is defined for sets of sentences only; the countermodel definition of
\cref{sec:levels-05} is restricted to sentences for the same reason. A
refutation of a consequence between open formulas states a structure and an
assignment.

\subsection{Assumptions, discharge, and derivability}
\label{sec:p01-w02-p01-016-assumptions}

Derivability in \cref{sec:levels-03} requires a derivation ``whose
undischarged assumptions are members of \(\Gamma\)''. For the
natural-deduction calculi on the handbook's route this is made exact as
follows. The rules of \(\mathsf{ND}_{\mathsf{CPL}}\) are displayed in
\cref{sec:prop-deduction-01}; \(\mathsf{ND}_{\mathsf{FOL}}\) adds the
quantifier, conversion and identity rules of \cref{sec:fol-deduction-01}.

\begin{hbdefinition}{def:p01-w02-p01-016-assumption}{Assumption, discharge, open and undischarged assumption}
Let \(D\) be a natural-deduction derivation: a finite rooted tree of formula
occurrences in which each inner node, with its immediate predecessors, is an
instance of a rule \cite[ch.~``Natural Deduction and Sequents'', definition
``Natural-deduction derivation'']{FrozenFormalLogicSemanticData}.
\begin{enumerate}
  \item An \emph{assumption occurrence} of \(D\) is a leaf that is not the
        conclusion of a rule without premises. In \(\mathsf{ND}_{\mathsf{CPL}}\)
        every leaf is an assumption occurrence; in
        \(\mathsf{ND}_{\mathsf{FOL}}\) the only premise-free rule is \(=I\),
        whose conclusion \(t=t\) is a leaf that is not an assumption.
        Assumption occurrences carry discharge labels drawn from
        \(\mathbb N\).
  \item A \emph{discharging rule} is one of \(\to I\), \(\neg I\),
        \(\mathsf{RAA}\), \(\lor E\) and \(\exists E\). An application of
        \(\to I\), \(\neg I\), \(\mathsf{RAA}\) or \(\exists E\) carries one
        label \(i\) and \emph{discharges} exactly the open occurrences marked
        \(i\) of its bracketed assumption (\(\varphi\), \(\varphi\),
        \(\neg\varphi\), \(\varphi[y/x]\) respectively) that lie above the
        premise to which the bracket belongs; an application of \(\lor E\)
        carries labels \(i,j\) and discharges the open occurrences of the
        first disjunct marked \(i\) above its second premise and those of the
        second disjunct marked \(j\) above its third premise. Each label
        carried by an application discharges at least one occurrence.
  \item An assumption occurrence is \emph{discharged} in \(D\) iff some rule
        application on the path from it to the root discharges it, and
        \emph{open}, or \emph{undischarged}, otherwise.
  \item \(\operatorname{Open}(D)\) is the set of formulas that have at least
        one open occurrence in \(D\).
\end{enumerate}
Then \(\Gamma\derives_{\mathcal C}\varphi\), for
\(\mathcal C\in\{\mathsf{ND}_{\mathsf{CPL}},\mathsf{ND}_{\mathsf{FOL}}\}\),
holds iff some \(\mathcal C\)-derivation \(D\) has conclusion \(\varphi\) and
\(\operatorname{Open}(D)\subseteq\Gamma\)
\cite[ch.~``Natural Deduction and Sequents'', definition ``Propositional
derivability''; ch.~``First-Order Deduction'', definition ``First-order
derivability'']{FrozenFormalLogicSemanticData}. The formula \(\varphi\) is a
theorem of \(\mathcal C\) iff \(\operatorname{Open}(D)=\varnothing\) for some
derivation \(D\) of \(\varphi\), that is, iff every assumption occurrence of
\(D\) is discharged.
\index{assumption!open}\index{assumption!undischarged}\index{discharge}
\end{hbdefinition}

Two features of the definition matter in use. Openness is a property of an
\emph{occurrence}: one formula may have a discharged occurrence and an open
one in the same derivation. And the side conditions of \(\forall I\) and
\(\exists E\) refer to the open assumptions of the subderivation at the point
where the rule is applied, so whether an assumption has already been
discharged decides whether a generalization is permitted.

\begin{hblemma}{thm:p01-w02-p01-016-open-bounds}{proved=annex:p01-w02-p01-016-open-bounds}{How a rule application changes the open assumptions}
\leavevmode
\begin{enumerate}
  \item If \(D\) consists of a single assumption occurrence of \(\varphi\),
        then \(\operatorname{Open}(D)=\{\varphi\}\); if \(D\) consists of a
        single conclusion of \(=I\), then \(\operatorname{Open}(D)=\varnothing\).
  \item Let \(D\) end with an application \(\rho\) of a rule to the
        conclusions of immediate subderivations \(D_1,\ldots,D_n\), with
        \(n\ge1\), and let \(\Delta_\rho\) be the set of formulas of which
        \(\rho\) discharges at least one occurrence. Then
        \[
          \Bigl(\bigcup_{k=1}^{n}\operatorname{Open}(D_k)\Bigr)
          \setminus\Delta_\rho
          \;\subseteq\;
          \operatorname{Open}(D)
          \;\subseteq\;
          \bigcup_{k=1}^{n}\operatorname{Open}(D_k),
        \]
        with equality throughout when \(\rho\) discharges nothing.
\end{enumerate}
\end{hblemma}

Both inclusions in clause 2 can be strict at the formula level, which is why
the open set is not a function of the open sets of the subderivations alone.
The example below exhibits a case where a formula in \(\Delta_\rho\) stays
open.

\begin{hbtheorem}{thm:p01-w02-p01-016-derivability-structure}{proved=annex:p01-w02-p01-016-derivability-structure}{Structural properties of derivability}
Let \(\mathcal C\) be \(\mathsf{ND}_{\mathsf{CPL}}\) or
\(\mathsf{ND}_{\mathsf{FOL}}\).
\begin{enumerate}
  \item If \(\varphi\in\Gamma\), then \(\Gamma\derives_{\mathcal C}\varphi\).
  \item If \(\Gamma\subseteq\Delta\) and \(\Gamma\derives_{\mathcal C}\varphi\),
        then \(\Delta\derives_{\mathcal C}\varphi\).
  \item If \(\Gamma\derives_{\mathcal C}\varphi\), then
        \(\Gamma_0\derives_{\mathcal C}\varphi\) for some finite
        \(\Gamma_0\subseteq\Gamma\).
  \item \(\derives_{\mathcal C}\varphi\) iff some derivation of \(\varphi\)
        has all its assumption occurrences discharged.
\end{enumerate}
\end{hbtheorem}

\begin{hbtheorem}{thm:p01-w02-p01-016-discharge}{proved=annex:p01-w02-p01-016-discharge}{Discharge theorem for natural deduction}
Let \(\mathcal C\) be \(\mathsf{ND}_{\mathsf{CPL}}\) or
\(\mathsf{ND}_{\mathsf{FOL}}\). For every set \(\Gamma\) and formulas
\(\varphi,\psi\),
\[
  \Gamma\cup\{\varphi\}\derives_{\mathcal C}\psi
  \iff
  \Gamma\derives_{\mathcal C}\varphi\to\psi .
\]
\end{hbtheorem}

\begin{hbexample}{def:p01-w02-p01-016-worked-discharge}{Tracking the open assumptions}
Let \(P,Q,S\in\mathsf{Prop}\). In the derivation below, written as a
numbered list in the inherited style, the last column records
\(\operatorname{Open}\) of the subderivation ending at each line.
\[
\begin{array}{@{}r@{\;\;}l@{\;\;}l@{\;\;}l@{}}
1.&[P]^1&\text{assumption}&\{P\}\\
2.&[P\to Q]^2&\text{assumption}&\{P\to Q\}\\
3.&Q&\to E,1,2&\{P,\,P\to Q\}\\
4.&[Q\to S]^3&\text{assumption}&\{Q\to S\}\\
5.&S&\to E,3,4&\{P,\,P\to Q,\,Q\to S\}\\
6.&P\to S&\to I,1\text{--}5&\{P\to Q,\,Q\to S\}\\
7.&(Q\to S)\to(P\to S)&\to I,4\text{--}6&\{P\to Q\}\\
8.&(P\to Q)\to((Q\to S)\to(P\to S))&\to I,2\text{--}7&\varnothing
\end{array}
\]
Line 6 witnesses \(\{P\to Q,Q\to S\}\derives_{\mathsf{ND}_{\mathsf{CPL}}}P\to S\);
line 8 shows that the final formula is a theorem. Each \(\to I\) removes one
formula from the open set, in agreement with
\cref{thm:p01-w02-p01-016-open-bounds}.

Now take two occurrences of \(P\) with different labels, \([P]^1\) and
\([P]^2\), combine them by \(\land I\) to \(P\land P\), and apply \(\to I\)
with label \(1\) to obtain \(P\to(P\land P)\). Only the occurrence marked
\(1\) is discharged. The occurrence marked \(2\) stays open, so this tree
witnesses \(\{P\}\derives_{\mathsf{ND}_{\mathsf{CPL}}}P\to(P\land P)\) and
not the theorem; for this tree, \(P\in\Delta_\rho\) for the final
\(\to I\), and still \(P\) belongs to its open set.
Marking both occurrences with one fresh label before applying \(\to I\), as
the inherited definition permits, discharges both and proves
\(\derives_{\mathsf{ND}_{\mathsf{CPL}}}P\to(P\land P)\).
\end{hbexample}

\begin{hbtheorem}{thm:p01-w02-p01-016-eigen-failure}{proved=annex:p01-w02-p01-016-eigen-failure}{An open assumption prevents generalization}
Let \(P\) be a unary relation symbol and \(x\) a variable.
\begin{enumerate}
  \item The two-node tree whose leaf is an assumption occurrence of
        \(P(x)\) and whose root \(\forall x\,P(x)\) is obtained by
        \(\forall I\) is not a derivation, and
        \(\{P(x)\}\not\entails\forall x\,P(x)\).
  \item \(\derives_{\mathsf{ND}_{\mathsf{FOL}}}\forall x\,(P(x)\to P(x))\).
\end{enumerate}
\end{hbtheorem}

In clause 2 the assumption \(P(x)\) is discharged by \(\to I\) before
\(\forall I\) is applied, so at the generalization step no open assumption
contains \(x\) free. The difference between the two clauses is exactly the
difference between an undischarged and a discharged assumption.

\begin{hbwarning}{def:p01-w02-p01-016-limits}{Limits of these definitions}
\begin{enumerate}
  \item The account of discharge applies to calculi with discharging rules.
        In the Hilbert-style modal calculi of
        \cref{sec:modal-temporal-epistemic-01}, a derivation from \(\Gamma\)
        is a finite sequence whose premise lines are never discharged, and
        necessitation applies only to formulas with a separate derivation
        from no premises \cite[ch.~``Modal, Temporal, and Epistemic Logic'',
        section ``Normal propositional modal logic'']{FrozenFormalLogicSemanticData}.
        There every premise used remains a premise, and
        \cref{thm:p01-w02-p01-016-discharge} is not available in the stated
        form.
  \item In a sequent calculus the role of the open assumptions of a
        derivation is played by the antecedent of its end sequent. The
        classical sequent section reads a finite multiset to the left of
        \(\derives\) through its support (\cref{sec:prop-deduction-02}), and
        the intuitionistic chapter defines the relation for an arbitrary set
        \(\Gamma\) through finite contexts drawn from \(\Gamma\)
        (\cref{ch:intuitionistic-logic}).
  \item \cref{thm:p01-w02-p01-016-unsat-reduction} assumes the classical
        negation clause. The intuitionistic forcing relation of
        \cref{ch:intuitionistic-logic} defines \(\neg A\) as
        \(A\to\bottom\), whose clause quantifies over all later worlds
        \cite[ch.~``Intuitionistic Logic'']{FrozenFormalLogicSemanticData},
        and the classical clause fails there: in the structure with worlds
        \(0\le1\) and \(V(p)=\{1\}\), the world \(0\) forces neither \(p\)
        nor \(\neg p\), since \(0\notin V(p)\) and the later world \(1\)
        forces \(p\). Clause 1 of that theorem is therefore not available for
        that semantics in the stated form.
\end{enumerate}
\end{hbwarning}

\paragraph{Relationships.}
\Cref{def:p01-w02-p01-016-datum,def:p01-w02-p01-016-satisfaction} make exact
the data and satisfaction relation named in the definition of a semantics in
\cref{sec:levels-01} and in the definition of semantic consequence in
\cref{sec:levels-03}; \cref{def:p01-w02-p01-016-assumption} makes exact the
undischarged assumptions in the definition of derivability there. Soundness
and completeness, defined in \cref{sec:levels-03}, connect the two relations
for each calculus and semantics; the first-order instances are
\cref{thm:fol-soundness,thm:fol-completeness}. The inherited propositional
completeness proof (\cref{thm:prop-completeness}) reaches consequence
through satisfiability: from \(\Gamma\not\derives\varphi\) it obtains that
\(\Gamma\cup\{\neg\varphi\}\) is consistent, constructs a valuation
satisfying it, and concludes \(\Gamma\not\entails\varphi\), which is the
semantic direction of clause 1 of
\cref{thm:p01-w02-p01-016-unsat-reduction}.

\paragraph{Sources.}
Every inherited definition used above is cited at its chapter and named
definition in the frozen manuscript
\cite{FrozenFormalLogicSemanticData}, whose text is reproduced on the
handbook's route at the cross-referenced sections. The results of this
section and their proofs in
\cref{sec:annex-p01-w02-p01-016-semantic-data} are editorial completions.
